% -----------------------------------------------------------------------------
% Purpose: Supplementary Information subsection for the murmurent manuscript,
%          on why approaches T_1-T_4 of the "Dance with Inhibition"
%          choreography are developed simultaneously rather than independently.
% Author:  bookworm agent, for @mhallet
% Date:    2026-07-28
% Input:   ~/repos/inhibition/decisions/D0001-D0026 (primary record);
%          ~/repos/murmurent_manuscript/main-article.tex (voice + labels, READ ONLY)
% Output:  a \subsection* to be pasted inside \begin{appendices}
%          \section{Supplementary Information}; cites entries in the companion
%          file manuscript_supplement_simultaneity_1.bib.
%
% NOTE: this file is NOT part of the manuscript repo. It is staged here so the
% manuscript can be pulled first, per rules/manuscript.md. Nothing in
% ~/repos/murmurent_manuscript was modified to produce it.
% -----------------------------------------------------------------------------

\subsection*{Why the four approaches are developed simultaneously}
\label{supp:simultaneity}

Expanding on \S\ref{results:choreography_pin1}, we set out why the four
approaches of the \emph{Dance with Inhibition} choreography are built
concurrently against one shared substrate rather than as four
independent projects that meet only at integration, and what that
choice costs.
The thesis is a trade and we state it as one: simultaneity buys
comparability and shared-control strength, and it spends error
independence.
Several of the choreography's design choices are best read as attempts
to keep the first while limiting the second.

\emph{Comparability is manufactured, not discovered.}
The four approaches answer one question with four instruments, so
integration is a comparison, and a comparison is only as good as the
axes it is made on; simultaneous development is what allows those axes
to be defined once.
$T_3$ and $T_4$ are both covalent, so they share a single covalent
protocol -- one pinned docking binary~\cite{mcnutt_gnina_2021}, one
version, one parameter set, one reactive-atom convention -- hashed into
a \texttt{protocol\_fingerprint} that both approaches carry with their
results.
The integration interface compares them only when those fingerprints
agree, and disables the within-covalent comparison rather than
displaying numbers produced under different rules; when the ligand form
changed (below), the fingerprint changed for both approaches together.
The descriptor, synthetic-accessibility and novelty axes are likewise
computed by one shared module for all four approaches, because the
interface pools them on one plot.
Developed separately, each approach would have arrived with a
defensible but private definition of ``novelty'' and the pooled plot
would have been meaningless.

\emph{Shared controls are built once and are stronger for it.}
Enrichment against property-matched decoys is expensive to do
well~\cite{mysinger_directory_2012,truchon_evaluating_2007}, and decoy
sets are known to carry biases that flatter methods which have learned
the bias rather than the
biology~\cite{chen_hidden_2019,sieg_need_2019,wallach_most_2018}, so
building it once on this receptor buys a better control than four
hurried ones.
That control then constrains more than one approach at a time: the
non-covalent gate (ROC-AUC $0.535$, $95\%$ CI $[0.215, 0.855]$,
EF$_{1\%}$ $0.0$) tells $T_1$ and $T_2$ together that their docking
rankings are weakly supported on this shallow, solvent-exposed
pocket~\cite{trott_autodock_2010,warren_critical_2006}, and the covalent
decoy set -- which we require to carry a warhead, since a decoy without
one is trivially separable -- serves $T_3$ and $T_4$ together.
Neither approach pair could have afforded that control alone.

\emph{Defects propagate, and so does their discovery.}
The strongest empirical argument for simultaneity in this build is that
a defect found in one approach was already latent in another.
The structural-alert screen excised the R-group from the scaffold and
capped the cut bond with hydrogen, which silently converts an amide into
a formamide and a thioether into a thiol; the alert
catalogue~\cite{brenk_lessons_2008} then correctly flags an aldehyde
that exists only in the fragment.
It surfaced because it rejected $97.7\%$ of $T_3$, whose generator
attaches decorations through carbonyls.
Repairing it also moved $T_4$'s gate ($1{,}683 \to 1{,}624$ passing).
$T_4$ alone would never have found it: its R-groups attach through
carbon and not one of its $1{,}782$ candidates carried the artefact.
The shared module is what made one approach's failure mode visible as
the other's silent one.

\emph{Seed orthogonality is engineered, and measured.}
The approaches are deliberately not equidistant from the seeds.
$T_2$ inherits ATRA's chemotype and $T_3$ and $T_4$ both inherit
sulfopin's, so $T_1$ is kept seed-free precisely so that agreement among
the others can be read as something other than shared ancestry.
Measured novelty (one minus maximum Tanimoto similarity to an external
reference set) has medians $0.87$ for $T_1$, $0.68$ for $T_3$ and $0.39$
for $T_2$, which is the intended ordering rather than an accident.
$T_3$ and $T_4$ are, in turn, a designed division of labour over one
seed: $T_4$ enumerates nine warhead classes against a fixed
$198$-member R-group library, $T_3$ fixes the warhead and spends its
whole budget on generative decoration.
Neither covers the two-dimensional warhead $\times$ R-group space alone;
the split is only coherent if the two are planned together.

\emph{The cost is correlated failure.}
Every argument above is an argument for shared infrastructure, and
shared infrastructure has a shared blast radius.
Two instances make this concrete.
Docking the pre-reaction ligand rather than the covalent adduct -- the
docking engine bonds a ligand atom to the receptor, it does not perform
reaction chemistry~\cite{scarpino_comparative_2018,bianco_covalent_2015,london_covalent_2014}
-- invalidated all $1{,}683$ of $T_4$'s covalent docks, and would have
invalidated $T_3$'s identically, because they share the protocol.
Earlier, an inappropriate protonation step dropped $28$ of $150$
receptor residues, Cys113 among them, while producing a file that still
looked like a protein and would have docked without error; one prepared
receptor serves all four approaches, so that single defect was capable
of destroying the entire choreography at once.
Independent implementations would have failed independently, and the
disagreement would itself have been a signal -- though the classic
result on multi-version redundancy is that even independently written
implementations of one specification do not fail
independently~\cite{knight_experimental_1986}, and here we have given up
even the appearance of independence in exchange for comparability.
The same source yields cheaper failures too: a stale reference to a
superseded warhead-class table caused $168$ of $192$ covalent docks to
fail, and shared staging plus contention for the same accelerators makes
the approaches schedule-coupled as well as method-coupled.

\emph{Coupled parameters force compromises on individual approaches.}
One protocol means neither covalent approach runs the protocol best
suited to itself.
The adduct-form transform is essential for $T_4$ and a no-op for $T_3$'s
acrylamide, yet $T_3$ carries the code path anyway -- deliberately,
since a second code path is exactly how two approaches drift apart.
Ranking within warhead class is necessary for $T_4$, whose classes dock
under different constraints and differ in size, and vacuous for
single-warhead $T_3$.
More seriously, a gate calibrated on one approach's chemistry is not
neutral toward another.
The covalent reactivity window is anchored on verified Pin1 covalent
actives that are all chloroacetamides, sulfamates or sulfonates -- one
chemical family -- and it excludes acrylamide, $T_3$'s expert-chosen
warhead, by $0.106$\,eV against a $0.5$\,eV tolerance we chose
ourselves.
Since computed electrophilicity is a weak proxy for measured labelling
kinetics~\cite{flanagan_chemical_2014}, the window flags rather than
vetoes, and flagged candidates proceed with their evidence attached.
That resolves the instance; it does not remove the structural hazard,
which is that shared evidence is never neutral between approaches that
occupy different regions of chemical space.

\emph{Convergence is weaker evidence than it looks.}
The deepest cost concerns the very thing simultaneity is supposed to
deliver.
Structural convergence -- a molecule surfaced by more than one approach
-- is offered as soft cross-validation needing no shared units.
But all four approaches dock into the same receptor under related
protocols and are annotated against the same reference set, so their
errors are correlated, and two approaches agreeing may report a shared
docking bias rather than a real signal.
This is the known limitation of consensus scoring, whose benefit derives
from the \emph{independence} of the combined methods: the classic
idealized analysis attributes consensus gains to the cancellation of
independent random error~\cite{wang_how_2001,charifson_consensus_1999},
and subsequent work finds that combining correlated scoring functions
delivers less than the independence assumption
implies~\cite{feher_consensus_2006,oda_comparison_2006,houston_consensus_2013}.
The same result appears in ensemble learning, where accuracy gains track
the diversity of the base
learners~\cite{dietterich_ensemble_2000,kuncheva_measures_2003}, and in
assay practice, where confirmation is sought from an orthogonal readout
rather than a replicate of the same
one~\cite{thorne_apparent_2010}.
Three things preserve part of convergence's value here.
$T_1$ is seed-free, so its agreement with a seeded approach is not
ancestral.
$T_3$ and $T_4$ propose a different mechanism from $T_1$ and $T_2$ --
covalent engagement of Cys113 rather than reversible occupancy -- so
agreement across that boundary crosses a real methodological divide.
And convergence is presented as an aid to a human reader rather than
folded into a forced consensus ranking; the four approaches report
incommensurable metrics and are displayed side by side with their gate
verdicts, which is the honest response to correlated evidence.

\emph{Does simultaneity harm the search?}
On the evidence of this build it does not harm the \emph{coverage} of
the search -- the seed allocation and the $T_3$/$T_4$ division of labour
widen it, and both are only possible under joint planning -- but it does
weaken the \emph{independence} of the search's internal checks.
Convergence between approaches is worth less here than the same
convergence would be worth between four separately implemented
pipelines, and we do not claim otherwise.
The mitigations are partial by construction: a seed-free approach and a
mechanistic split reduce correlation without eliminating it, and no
arrangement of a shared receptor and a shared scoring function makes
four approaches independent.
The choreography trades a weaker cross-validation signal for a stronger
set of shared controls and a genuinely comparable integration, and the
shared controls are the part we can measure.
